\clearpage\section{MOS 電晶體電路圖}
來源：\texttt{01\_dc\_transfer.sp} 與 \texttt{02\_noise\_cascade.sp}。三個並列支路 $k=0,1,2$ 的 $W_p$ 分別為 $0.45,1.8,7.2\,\mu$m；皆使用 $W_n=0.6\,\mu$m、$L=0.25\,\mu$m。DC 網表只含第一級 Xk；雜訊暫態網表在各支路加入同尺寸的 XBUFk。

圖中沿用所列網表的元件與節點名稱；相同節點標籤表示電氣相連。除另有註明外，NMOS 的 bulk 接地，PMOS 的 bulk 接 $V_{DD}$。三端符號省略 bulk 導線以保持圖面清楚。
\begin{figure}[!htbp]\centering
\resizebox{12cm}{!}{\begin{circuitikz}[american,font=\small]

\begin{scope}[shift={(0,0)}]
\node[pmos] (Xkp) at (0,1) {};
\node[nmos] (Xkn) at (0,-1) {};
\draw (Xkp.S) -- (0,2) node[above] {$V_{DD}$};
\draw (Xkn.S) -- (0,-2) node[ground] {};
\draw (Xkp.D) -- (0,0) -- (Xkn.D);
\draw (Xkp.G) -- (-1.15,1) -- (-1.15,-1) -- (Xkn.G);
\draw (-1.15,0) node[circ] {} -- (-1.8,0) node[above] {in};
\draw (0,0) node[circ] {} -- (0.85,0) node[above] {$y_k$};
\node[right] at (0.25,1) {MP};
\node[right] at (0.25,-1) {MN};
\node at (0,2.6) {};
\coordinate (Xkin) at (-1.8,0);
\coordinate (Xkout) at (0.85,0);
\end{scope}

\begin{scope}[shift={(5,0)}]
\node[pmos] (XBUFkp) at (0,1) {};
\node[nmos] (XBUFkn) at (0,-1) {};
\draw (XBUFkp.S) -- (0,2) node[above] {$V_{DD}$};
\draw (XBUFkn.S) -- (0,-2) node[ground] {};
\draw (XBUFkp.D) -- (0,0) -- (XBUFkn.D);
\draw (XBUFkp.G) -- (-1.15,1) -- (-1.15,-1) -- (XBUFkn.G);
\draw (-1.15,0) node[circ] {} -- (-1.8,0) node[above] {$y_k$};
\draw (0,0) node[circ] {} -- (0.85,0) node[above] {$z_k$};
\node[right] at (0.25,1) {MP};
\node[right] at (0.25,-1) {MN};
\node at (0,2.6) {};
\coordinate (XBUFkin) at (-1.8,0);
\coordinate (XBUFkout) at (0.85,0);
\end{scope}
\draw (Xkout) -- (XBUFkin);
\end{circuitikz}}
\caption{一個完整四顆 MOS 雜訊串接支路；三種寬度比各自獨立重複此接法。共用輸入並不代表各支路輸出相連。}
\end{figure}
